\subsection{Sprint 1}

Para esta sprint, estavam previstos três épicos, um para cada squad: Tags, Gerenciamento de Eventos e Home/Feed. Por natureza, nem todos os épicos eram iguais. Meu squad ficou responsável pelo épico de Tags, que era o menor; por isso, também nos comprometemos a ajudar os outros squads mais para o final da sprint. Dessa forma, acabamos atuando livremente em tasks de outras equipes.

Da minha parte, fiquei muito focado na comunicação com os membros do squad, principalmente os mais iniciantes, para garantir que estivessem engajados e entendessem bem o projeto. Ainda assim, fiz questão de não me limitar ao meu squad, apenas dando mais atenção a ele.

O time conseguiu entregar tudo o que estava previsto na sprint e, em nenhum momento, tivemos o desespero de estar fora do cronograma. Alguns detalhes precisaram ser implementados após o code freeze, mas, em geral, o trabalho da equipe foi bom. Isso é ainda mais relevante considerando que os Ages III mantiveram a baseline da qualidade do código mergeado muito alta, exigindo diversas mudanças nas branches antes do merge.

Sinto que também consegui fazer um bom trabalho na gestão, garantindo o engajamento da equipe e mantendo a comunicação com todos os membros do squad. Mantive contato com outros membros do time que não fazem parte do meu squad, principalmente Ages 3 e Ages 4 com quem tenho cadeiras, o que me permitiu acompanhar a evolução da equipe com mais facilidade. Além disso, entreguei algumas tasks mais urgentes para garantir que outras fossem desbloqueadas, buscando liderar pelo exemplo e aliviar um pouco a carga dos Ages 1 e Ages 2, principalmente neste começo de Ages, no qual ainda estão se familiarizando com o novo ambiente.

Mesmo com todos os pontos positivos relatados e com a qualidade da entrega, houve alguns pontos negativos. Terminamos a retrospectiva com três action items. O primeiro foi o rebalanceamento dos squads. Isso, por si só, não seria um problema, mas uma integrante Ages 1 de um dos squads desistiu no meio do projeto por não ter sido integrada pela equipe. Em segundo lugar, refletimos que, para manter as correções de PRs tão estritas quanto estavam, precisaríamos melhorar muito a comunicação. Um PR que recebesse pedidos de alteração em três momentos diferentes poderia demorar mais de uma semana para ser mergeado e, tratando-se de componentes, isso praticamente inviabiliza nosso desenvolvimento enquanto equipe da Ages. Por último, notamos também uma baixa adesão às dailys de sábado. Tentamos ser o mais flexíveis possível, permitindo daily por mensagem, enviada tanto no sábado quanto no domingo. Como na Ages passada tivemos um problema semelhante de adesão, essa flexibilização não foi suficiente para solucioná-lo. Podemos reforçar e melhorar a comunicação individual nos squads para que a falta de adesão não prejudique o planejamento e a gerência do projeto.

Uma das primeiras lições aprendidas foi relacionada à correção de PRs. Percebi que uma correção muito estrita atrasa o ciclo de desenvolvimento, pois as branches demoram muito para ser mergeadas. Conversei com os Ages 3 para adaptar um pouco esse processo em componentes e tarefas bloqueantes, tendo em vista que alguns componentes do frontend ficaram travados por mais de uma semana, inviabilizando o começo de tasks de telas que dependiam deles. A dificuldade é encontrar um equilíbrio entre desbloquear as próximas tarefas o mais rápido possível e, ao mesmo tempo, garantir a manutenção da qualidade do código. Além disso, como relatado anteriormente, tivemos um aprendizado relativo às dailys de sábado: minha hipótese da Ages passada foi refutada, tendo em vista que tivemos um problema de adesão semelhante mesmo tornando o processo menos estrito.

A próxima sprint terá duas semanas e, provavelmente, um code freeze na segunda. Por isso, planejamos um período mais corrido do que o da sprint 1. Por outro lado, acredito que o time tenha evoluído e consiga abraçar o mesmo número de tasks da última sprint com menos tempo de desenvolvimento. Planejo, junto com os membros do time, entregar logo os componentes do frontend para desbloquear todas as telas, visto que, nesta sprint, ficamos muito bloqueados por isso e, por sorte, tínhamos mais dias. Não devemos cometer esse erro novamente, ou fracassaremos na entrega da próxima sprint. Além disso, planejo continuar oferecendo suporte aos membros da equipe, ajudando desde os Ages 1 até outros Ages 4 que precisarem e procurando ter sempre certeza de que todos estejam engajados.
